\section{Data and Evaluation Protocol}
\label{sec:data-evaluation}

\subsection{xView3 Study Cohort}

We use xView3-SAR Sentinel-1 ground-range-detected scenes with VH and VV
polarizations at 10~m pixel spacing~\cite{paolo2022xview3}. Scene identifiers,
not image chips, are the unit of partitioning. A frozen seed-0 cohort contains
150 study scenes: 111 train, 23 development, and 16 test. Selection is
stratified by coarse geographic region and shoreline presence. The split is
immutable, and every data loader checks that scene membership is disjoint.
The exact counts and evaluation supports are summarized in
Table~\ref{tab:data-protocol}.

\begin{table}[t]
\caption{Frozen scene-level data protocol and annotation support under the
corrected ground-truth contract. The 50 verified scenes remain sealed until
the held-out test matrix and review gates complete.}
\label{tab:data-protocol}
\centering
\small
\setlength{\tabcolsep}{7pt}
\begin{tabular}{@{}lrrrr@{}}
\toprule
Partition & Scenes & Vessel H/M & Non-vessel H/M & LOW ignore \\
\midrule
Train (full) & 111 & -- & -- & -- \\
Training-time dev subset & 8 & 517 & 107 & 118 \\
Development & 23 & 1,479 & 804 & 441 \\
Test & 16 & 1,165 & 420 & 325 \\
Verified final & 50 & \multicolumn{3}{c}{sealed; human verified} \\
\bottomrule
\end{tabular}
\vspace{2pt}
\parbox{\textwidth}{\small Nested training subsets contain 12, 28, 56, and
111 scenes for the nominal 10\%, 25\%, 50\%, and 100\% fractions.}
\end{table}


The 111 training scenes produce 105,408 overlapping $800\times800$ chips.
Training-time normalization statistics are computed from training scenes
only: VH mean/standard deviation $-26.448/5.951$~dB and VV
$-16.599/6.062$~dB. Each optimization sample is a shared $512\times512$ crop
from these chips. The label-fraction subsets contain exactly 12, 28, 56, and
111 scenes and are nested as in Eq.~\ref{eq:nesting}; no development, test, or
human-verified scene appears in any training subset.

A separate set of 50 human-verified scenes is designated
\texttt{eval\_final}. It is disjoint from the 150-scene study cohort and is
used once, only after the 32-cell cohort, model selection, operating
thresholds, and analysis code are frozen. Human verification supplies the
manual-label provenance needed for the dark-vessel slice. Consequently,
dark-vessel recall is undefined for the ordinary development and test sets and
is reported only on \texttt{eval\_final}.

\subsection{Ground-Truth Contract}

Evaluation uses the same target semantics as training. A row is a positive
vessel point only when \texttt{is\_vessel} is true and confidence is HIGH or
MEDIUM. A HIGH/MEDIUM row explicitly marked as non-vessel is background and is
excluded from the positive set. LOW-confidence rows are retained as ignore
points. Confidence and boolean fields are normalized and validated before
constructing scorer inputs; an ambiguous HIGH/MEDIUM vessel flag is an error,
not an implicit positive.

This contract yields 517 positive vessels, 107 excluded HIGH/MEDIUM
non-vessels, and 118 LOW-confidence ignores on the fixed eight-scene
training-time development subset. The full 23-scene development set contains
1,479/804/441 rows in those three categories, and the 16-scene test set
contains 1,165/420/325. These support counts are pre-registered integrity
checks: evaluation stops if regenerated counts differ. Training jobs can read
only the 111 training scenes and fixed development eight; the test view is
staged only after an immutable 32-cell cohort exists, and final-evaluation
labels remain outside both views.

An audit of an earlier V100 diagnostic campaign found that its development
conversion had admitted HIGH/MEDIUM non-vessels as positives. The same audit
found that held-out scripts could pair the selected \texttt{best.ckpt} with a
threshold from the last development evaluation rather than the threshold
selected for that checkpoint. Because development F1 controls early stopping
and checkpoint selection, post-hoc rescoring cannot make those runs
reportable. We corrected the caller-side ground-truth conversion without
modifying the frozen scorer, bound each best-development operating point to an
exact checkpoint hash and epoch, and require all 32 H100 cells to restart from
scratch. V100 values are excluded from reportable tables, curves, summaries,
and cross-hardware aggregates.

\subsection{Checkpoint Selection and Held-Out Access}

Every five epochs, each training run performs full-scene inference on the same
fixed, sorted set of eight development scenes. Candidate detections are swept
to select the score threshold that maximizes aggregate development F1. The
complete operating point---epoch, threshold, precision, recall, counts, and
candidate total---is stored together with the SHA-256 of the exact best
checkpoint. Resume and requeue restore this state. A completed cell is invalid
if checkpoint epoch, checkpoint digest, or operating-point record disagree.

After all 32 corrected cells form one immutable cohort, the stored threshold
is applied verbatim to the 16 test scenes. Test annotations are never used to
retune a threshold, select a checkpoint, or change an analysis choice. The
same checkpoint-bound threshold is subsequently used for the once-only
human-verified evaluation. This separation makes the development set the sole
source of model selection and operating-point calibration.

\subsection{Metrics and Reporting}

The frozen scorer performs one-to-one geographic matching within 200~m.
Aggregate precision, recall, and F1 are computed from sums of true positives,
false positives, and false negatives across scenes; the principal
label-efficiency outcome is aggregate test F1 at the development-selected
threshold. The scorer implementation and hash are held fixed across every
arm.

Two operational slices are evaluated on the human-verified set. Dark-vessel
recall measures matched manually sourced vessel positives, which serve as the
benchmark's proxy for vessels without matched AIS. We report recall, not a
dark-vessel F1, because ``dark'' is a ground-truth attribute and is not defined
for arbitrary false-positive predictions. Near-shore F1 uses vessel ground
truth and unmatched predictions at or within 2~km of shore, so coastal false
alarms count against the slice. The ordinary test split contains only two
valid near-shore vessel positives; its near-shore F1 is therefore identified
as a sparse diagnostic, while the registered operational near-shore result
comes from the larger human-verified set.

The main results report all eight four-point H100 curves, the transfer gains
in Eq.~\ref{eq:transfer-gain}, the SAR--optical contrasts in
Eq.~\ref{eq:sar-optical-gap}, measured epochs and GPU-hours, and the once-only
verified metrics. We neither enforce monotonicity nor select favorable
fractions post hoc. With one seed, these are point estimates without error
bars, confidence intervals derived from seed variation, or significance
claims.

Finally, the frozen scene split is leakage-free but not a claim of geographic
independence. A post-freeze proximity audit found that 22 of 23 development
scenes lie within 32~km of a training scene and 42 of 50 human-verified scenes
lie within 50~km of one. We therefore describe held-out performance as
\emph{in-region generalization}; conclusions about transfer to new maritime
regions require a separately designed geographic holdout.
